\section{Limiting bounds and primitive FBBT updates}\label{app:fbbt}
Feasibility-based bound tightening (FBBT) contracts a box by propagating
individual constraints. Its limit, the number of primitive updates, and
the finite tolerance used by an implementation are distinct objects.
When the limiting box is nonempty, continuous linear FBBT limits can be
obtained by an LP \citep[Theorem~4.1]{BelottiEtAl2012}, although primitive
iteration need not terminate finitely.
Here all variables start in $[0,1]$. A \emph{primitive update} computes the
interval hull for one affine equality or one nonnegative product equality
and intersects it with the current box. For the limiting-bound theorem
we allow any additional sound contracting updates, provided every defining equation
receives its forward natural interval-arithmetic update infinitely often.
This is fairness, with no bounded-delay assumption. Exact real arithmetic
specifies the mathematical limit; an approximation algorithm is measured
in the ordinary binary Turing model.

\begin{lemma}[Least fixed point transfer]\label{lem:fbbt-fixed-point}
Let $P$ have nonnegative polynomial coefficients, and suppose $x=P(x)$
has a solution in the unit cube. Under sound contracting updates and fair forward
propagation, its limiting lower vector is the least nonnegative fixed
point $q$ of $P$. This includes lifted arithmetic-node equations.
\end{lemma}
\begin{proof}
The sequence $s^0=0$, $s^{j+1}=P(s^j)$ increases and is bounded by every
nonnegative fixed point. Its limit $q$ is a fixed point by continuity and
is the least one. Soundness preserves $q$, so every lower vector $l$
satisfies $l\le q$. The natural interval lower endpoint is exactly $P_i(l)$.
Inductively, once $l\ge s^j$, fairness supplies a finite later time at
which all coordinates have received a forward update and $l\ge s^{j+1}$.
Hence the lower limit is at least every $s^j$ and equals $q$. Apply the
same argument to the augmented polynomial system for an expression graph.
\end{proof}

\begin{proposition}[Restricted limiting-bound hardness]\label{prop:fbbt-hardness}
Computing a designated limiting lower bound to additive error $1/4$ is
PosSLP-hard, even for feasible affine and bilinear defining systems in
the unit cube with at most three variable names per equation,
nonnegative coefficients in $\{1/2,1\}$, constants in $\{0,1/2,1\}$,
and distinct names in every product. Every strongly connected component
of the directed defining-equation dependency graph has at most four
vertices. The feasible set has at most two points, all rational.
\end{proposition}
\begin{proof}
PosSLP asks whether an integer arithmetic circuit using $0,1,+,-,\times$
has positive output. Following the circuit normalization and amplification
mechanism of \citet[Theorem 5.2]{EtessamiYannakakis2009}, replace every
gate by positive and negative parts. Multiplication uses
$(uv)^+=u^+v^++u^-v^-$ and $(uv)^-=u^+v^-+u^-v^+$; addition and
subtraction use the corresponding sums. It suffices to compare two
nonnegative integer outputs $U,V$.

Pad these fan-in-two circuits to common depth $L\ge1$ with alternating
addition and multiplication layers, using identities $v+0,v\cdot1$.
This costs polynomial size. Define proof-only normalizers $M_0=1$,
$M_j=2M_{j-1}$ on addition layers and $M_j=M_{j-1}^2$ on multiplication
layers. Encode each gate by its normalized value $p$ and complement $r$.
At inputs they are zero or one; the equations for subsequent gates are
\[
 p=(p_j+p_k)/2,\quad r=(r_j+r_k)/2
\]
for addition and
\[
 p=p_jp_k,\quad v=p_jr_k,\quad r=r_j+v
\]
for multiplication. These identities give $p=\operatorname{value}/M_j$
and $r=1-p$, with unique values in $[0,1]$. Introduce a copy variable
whenever a product would repeat a name. Each equation has at most three
names and all these dependencies are acyclic. With $M=M_L\le2^{2^L}$,
put
\[
 c=(p_U+r_V)/2=\tfrac12+(U-V)/(2M),\qquad
 d=(r_U+p_V)/2=1-c.
\]
The four-equation component
\[
 a'=a,\quad t=aa',\quad v=ct,\quad a=d+v
\]
implements $a=d+ca^2$. Its least nonnegative root $a_*$ equals one if
$c\le1/2$ and $d/c$ if $c>1/2$. In the latter case, writing
$\Delta=U-V\ge1$,
$1-a_*=2\Delta/(M+\Delta)\ge1/M$.

Construct $b=2^{-2^{L+2}}$ and its complement by $L+2$ paired squarings
starting from $1/2$. These numbers require polynomially many equations,
not their full rational expansions as input constants. They satisfy
$b\le1/(8M)$. Add $e=(1-b)a$, represented using the complement variable,
and the two-equation component $w=ez$, $z=b+w$.
At the least detector root it has the unique solution
\[
 z_*=\frac{b}{1-(1-b)a_*}
 =1\quad(U\le V),\qquad z_*\le bM\le1/8\quad(U>V).
\]
All displayed values lie in the cube and form a fixed point. The least
fixed point of the full system has the unique acyclic values, the least
detector root (bounded above by this displayed fixed point), and then
this unique final value. Lemma~\ref{lem:fbbt-fixed-point} identifies $z_*$
with the FBBT limit. An additive-$1/4$ answer is at least $3/4$ in the
first case and at most $3/8$ in the second; threshold $1/2$ decides PosSLP.
The largest directed component has the four detector variables, and the
last has two. The detector has at most two rational roots, and each
admissible root determines all other variables uniquely and rationally.
\end{proof}
The directed graph has an edge from a defined variable to each argument
on its right side; this is not an undirected treewidth restriction.
PosSLP is not known to be NP-hard. The proposition does not assert
hardness of executing a prescribed finite number of sweeps or of obtaining
an approximately stationary box. A unique rational feasible point is
possible if one further affine inequality is allowed: introduce $y=ca$
and impose $y\le1/2$. It retains the least detector root and excludes
$a=1$ when $c>1/2$; soundness and the fair original forward updates retain
the same limiting-bound proof.

\begin{proposition}[Every primitive schedule can be slow]\label{prop:fbbt-iterations}
For each $n\ge1$ a constant-coefficient affine--bilinear system with
$O(n)$ variables and equations requires at least $2^{2^n-1}$ applications
of one affine constraint before a designated lower bound reaches $1/2$,
although its limit is one. This holds for every fair ordering of the
bidirectional primitive interval-hull updates defined above. The only
feedback component has two variables and is linear after its upstream
values are fixed.
\end{proposition}
\begin{proof}
Start with $b_0=c_0=1/2$ and, for $i=1,\ldots,n$, impose
\[
 h_i=b_{i-1},\quad b_i=b_{i-1}h_i,\quad
 v_i=b_{i-1}c_{i-1},\quad c_i=c_{i-1}+v_i.
\]
These acyclic equations give $b=b_n=2^{-2^n}$, $c=c_n=1-b$.
Add $w=cz$, $z=b+w$. The unique feasible values are $z=1,w=c$.
For a stronger initialization fix all acyclic variables exactly, leaving
$z,w\in[0,1]$. Primitive interval-hull contractors are monotone under box
inclusion, so the original run with the same schedule has no larger lower
bounds.

In the stronger run the invariant is $0\le l_w\le cl_z$, $l_z\le1$.
A product update sets $l_w$ to $cl_z$ at most; its inverse bound
$l_z\ge l_w/c$ cannot improve $l_z$. An affine update sets
\[
 l_z^{\rm new}=\max(l_z,b+l_w)\le b+cl_z,\qquad
 l_w^{\rm new}=\max(l_w,l_z-b).
\]
Since $l_z-b-cl_z=b(l_z-1)\le0$, both terms defining the new lower bound
on $w$ are at most $cl_z^{\rm new}$. This proves the invariant, including
simultaneous hull updates. Upper endpoints cannot improve these formulas:
feasibility forces $u_z=1$, while $u_w$ is either one or $c$.
After $K$ affine updates, induction and Bernoulli's inequality give
\[
 l_z\le1-c^K=1-(1-b)^K\le Kb.
\]
Thus $l_z\ge1/2$ requires $K\ge1/(2b)=2^{2^n-1}$.
\end{proof}
The precise bound is in the displayed parameter $n$; encoding indexed
variable names as binary strings gives ordinary input length $O(n\log n)$.
At $l_z=l_w=0$ with upstream variables fixed, every primitive update changes
any lower endpoint by at most $b$, although $l_z$ is distance one from its
limit. Small local residual and small distance to the limiting box are
therefore different guarantees.

Slow and nonfinite linear FBBT iteration is classical, as are repeated
squaring and error amplification. The slow-Newton examples of
\citet[Section 7, Theorem 7.1]{EsparzaEtAl2010} and
\citet[Section 4.1, equation (18)]{StewartEtAl2015} also imply a doubly
exponential first-bit bound for ordinary Kleene iteration. This is an
inference from their displayed systems, not their stated Newton theorem.
For Stewart et al.'s chain, $x_0=(x_0^2+1)/2$ and
$x_i=(x_i^2+x_{i-1})/2$, the errors from the fixed point one satisfy
$e_0(k+1)=e_0(k)-e_0(k)^2/2$ and $e_i(k)^2\ge e_{i-1}(k)$, the latter
because each iterate is below its next update. The first recurrence gives
$e_0(k)\ge1/(k+1)$ by induction (the map is increasing on $[0,1]$).
Hence $e_n(k)\ge(k+1)^{-1/2^n}>1/2$ when $k+1<2^{2^n}$. The
specific scope proved here is every ordering of the stated bidirectional
primitive contractors, with one feedback component that becomes linear.
Symbolic substitution, equation aggregation, fixed-point acceleration,
and global contractors are outside this iteration bound.
